%% Elements of Nested Fibrational Cosmology
%% BIO Branch Book: Biology Branch (Prospective)
%% Status: Architecturally constituted; pre-closure — all theorem shells carry open obligations
%%
%% Status legend:
%%   [U]  Unconditional  -- proved from spine imports and prior [U] results
%%   [C]  Conditional    -- depends on explicitly declared hypotheses
%%   [D]  Definition-structural
%%   [R]  Remark only
%%   [O]  Open obligation / named failure frontier
%%   [B]  Bridge (Book VI licensed continuum surrogate)

\documentclass[12pt,oneside]{amsbook}
\usepackage{amsmath,amssymb,amsthm}
\usepackage{mathrsfs}
\usepackage{enumitem}
\usepackage{hyperref}
\usepackage{geometry}
\geometry{margin=1.2in}

\theoremstyle{plain}
\newtheorem{theorem}{Theorem}[section]
\newtheorem{lemma}[theorem]{Lemma}
\newtheorem{proposition}[theorem]{Proposition}
\newtheorem{corollary}[theorem]{Corollary}

\theoremstyle{definition}
\newtheorem{definition}[theorem]{Definition}
\newtheorem{standingrule}[theorem]{Standing Rule}
\newtheorem{openobligation}[theorem]{Open Obligation}

\theoremstyle{remark}
\newtheorem{remark}[theorem]{Remark}
\newtheorem{scholium}[theorem]{Scholium}

\newcommand{\status}[1]{\textup{[#1]}}

%% Notation
%%   R_Bio               Biology target regime
%%   O_Bio               Biology observable family = (Phi_Bio, M_Bio, H_Bio)
%%   Phi_Bio             Replication kernel (quotient-visible self-copy data)
%%   M_Bio               Metabolic ledger (resource-transformation bookkeeping)
%%   H_Bio               Heredity-descent object (heritable transport map)
%%   F_Bio               Fidelity margin (replication fidelity separation quantity)
%%   B_Bio               Boundary-capacity object (membrane/interface structure)
%%   E_Bio               Biology endpoint datum
%%   Rep(X; t)           Replication signature of configuration X at step t

\begin{document}

\title{BIO Branch Book\\[6pt]
 Replication, Metabolism, and Certified Biological Structure}
\author{Elements of Nested Fibrational Cosmology\\
 Prospective Derived Branch}
\date{}
\maketitle
\tableofcontents
\newpage

%% ---------------------------------------------------------------
\section{Purpose and Canonical Seed}\label{sec:bio-purpose}
%% ---------------------------------------------------------------

This branch asks whether biological structure --- self-replication,
metabolism, heredity, cellular boundary formation, and Darwinian
evolution --- can be realized as a lawful descendant of the NFC source
architecture. The aim is not to import the standard language of
cells, genes, fitness landscapes, proteins, or metabolic pathways as
primitives, but to determine when such language is licensed
as an effective encoding of certified discrete observable structure
within a declared branch regime. This posture is mandated by
Book~I's no-smuggling stance, Book~V's branch-legitimacy law,
Book~VI's continuum-interface discipline, and Book~VII's
scoped-certificate law.

The central question is: starting from nothing but a finite class of
physical configurations, a finite admissible probe discipline, and the
observational equivalence imposed by that discipline, can one derive
certified self-replicating structure, metabolic bookkeeping, heritable
transport, and the emergence of a boundary-enclosed arena? No cell,
gene, fitness landscape, vitalistic force, or biological essence is
assumed. Each such notion,
if it appears at all, must emerge as a quotient-visible,
source-descended object within the declared regime.

Three structural observations motivate the branch within the existing
canonical corpus. First, Book~II already establishes a
\emph{boundary capacity law} --- information capacity at interfaces is
logarithmically bounded. The biology branch will ask whether a
certified observable configuration can develop a \emph{self-organized
boundary} that exploits this capacity structure without violating it.
Second, Book~III already establishes a \emph{unified coercive-transport
inequality} governing how stability costs distribute. The biology
branch will ask whether a configuration can minimize that transport
cost by self-replicating --- reproducing its own certified observable
class rather than diffusing into equivalence-class noise. Third, the
NS branch already treats a \emph{continuation criterion} whose failure
marks the singular endpoint; the biology branch treats the analogous
quantity --- a \emph{replication-continuation criterion} whose
maintenance marks the persistence of a living configuration.

The initial endpoint is deliberately modest: certified replication
structure, metabolic bookkeeping, and hereditary transport within a
declared observable regime. Multicellular organization, evolutionary
dynamics, and molecular-biological specificity are deferred as named
open obligations.

\begin{remark}[\status{R}]\label{rem:bio-strategic-path}
\textbf{Strategic path.}
The shortest canonical route is replication-founded heredity: begin
from certified admissible-copy events (the replication kernel), build
a metabolic ledger for resource-transformation bookkeeping, show that
replication is fidelity-bounded by a positive margin, prove a
no-hidden-channel theorem for the hereditary transport map, then ask
whether the assembled structure certifiably descends from $\mathbf{U}$
via the Book~V branch-legitimacy law. Evolutionary dynamics and
multicellular organization remain downstream open obligations.
\end{remark}

%% ---------------------------------------------------------------
\section{Imported Spine Structure}\label{sec:bio-imports}
%% ---------------------------------------------------------------

\textbf{Imported spine results.} All certified results of Books~I--VII
are required for lawful descent, transport accounting, continuum-interface
discipline, and closure governance. In particular:
\begin{enumerate}[label=\textup{(\arabic*)}]
 \item \emph{Book~I}: observational quotient discipline; no primitive
 cells, genes, organisms, species, or molecular-biological
 vocabulary may be assumed. Biological content must be
 quotient-visible and admissibly testable, not imported as
 background life-science.
 \item \emph{Book~II}: stabilized local outcomes, collar-visible
 control, and boundary capacity structure. The logarithmic
 boundary capacity law is the canonical constraint on any
 self-organized interface claiming to separate an interior from an
 exterior without hidden information import.
 \item \emph{Book~III}: persistence and transport law; additive
 bookkeeping. The unified coercive-transport inequality governs how
 stability costs distribute across replication, metabolic
 transformation, and boundary maintenance --- the three canonical
 cost categories of the biology branch.
 \item \emph{Book~IV}: the universal source object $\mathbf{U}$ and
 source-descended comparison architecture. All certified biological
 objects must factor through $\mathbf{U}$.
 \item \emph{Book~V}: lawful branchhood and branch visibility; the
 constitutional law under which a biology branch is admitted.
 \item \emph{Book~VI}: continuum language licensed only as faithful
 shorthand for certified discrete structure within a declared
 interface regime; the main gate for continuous-time population
 dynamics, thermodynamic quantities, information-rate quantities,
 and chemical-kinetics notation in formal biological models.
 \item \emph{Book~VII}: scoped certificates, promotion law, named
 frontier stability, and prohibition of hidden upgrade. Claims
 about universal features of life, infinite evolutionary potential,
 or species-level universals are subject to this governance without
 exception.
\end{enumerate}

\textbf{Imported branch precedents.}
\begin{enumerate}[label=\textup{(\arabic*)}]
 \item \emph{Book~II boundary capacity}: the interface constraint that
 bounds the information capacity of any self-organized boundary;
 the \emph{membrane template} for this branch.
 \item \emph{NS branch}: the continuation-criterion / obstruction-decay
 pattern as the backbone of a main frontier theorem; the
 \emph{replication-continuation template} --- a biological
 configuration either continues to replicate or encounters a
 frontier analogous to NS blowup.
 \item \emph{SPEC branch}: probe-response objects and transition
 ledgers; the \emph{metabolic-probe template} for certified
 resource-transformation bookkeeping.
 \item \emph{LING branch}: the no-smuggling posture on inherited
 structure; in particular, the principle that a derived inventory
 (contrast classes, certified copy classes) must emerge from
 quotient structure, not be imported as a named alphabet or genetic
 code.
\end{enumerate}

%% ---------------------------------------------------------------
\section{Branch-Specific Arena}\label{sec:bio-arena}
%% ---------------------------------------------------------------

\subsection{Standing Rules}\label{sec:bio-standing-rules}

\begin{standingrule}[\status{D}]\label{sr:bio-no-smuggling}
\textup{(No Smuggling.)}
No biological object --- cell, gene, organism, chromosome, protein,
metabolic pathway, fitness value, or species --- may be treated as
primitive branch data merely because it is familiar in external
biological theory. Every branch-visible biological object must be
sourced through $\mathbf{U}$, the declared descent machinery, and the
biology observable family. This is a branch-local instance of
Book~I no-smuggling and Book~V branch-legitimacy discipline.
\end{standingrule}

\begin{standingrule}[\status{D}]\label{sr:bio-no-vitalism}
\textup{(No Vitalism.)}
No vitalistic force, teleological drive, biological essence, or
special-purpose self-organization principle may be assumed. Every
instance of self-organized structure that appears in this branch must
be derived from quotient-visible observational data and the spine's
transport and boundary machinery.
\end{standingrule}

\begin{standingrule}[\status{D}]\label{sr:bio-quotient-visible}
\textup{(Quotient Visibility.)}
A biological distinction between two configurations counts only when
it is quotient-visible under admissible probe tests. In particular,
a distinction between a living and a non-living configuration counts
only when it is witnessable by some admissible probe within the
declared regime --- not by appeal to an external biological authority.
\end{standingrule}

\begin{standingrule}[\status{D}]\label{sr:bio-no-hidden-gene}
\textup{(No Hidden Genetic Code.)}
No hidden genetic code, underlying molecular representation, or
abstract hereditary level may be assumed. Hereditary structure may
appear only after a branch-local derivation certifies that such
structure faithfully encodes quotient-visible replication data within
the declared regime.
\end{standingrule}

\begin{standingrule}[\status{D}]\label{sr:bio-no-self-promotion}
\textup{(No Self-Promotion.)}
A finite replication certificate may not self-promote to a claim about
unlimited evolutionary potential, infinite hereditary lineages, or
universal features of life. Such promotion requires an explicit
transfer theorem. This is a direct Book~VII governance consequence.
\end{standingrule}

\begin{standingrule}[\status{D}]\label{sr:bio-frontier-visibility}
\textup{(Frontier Visibility.)}
Named frontiers in replication fidelity, metabolic closure, hereditary
transport, boundary self-organization, evolutionary dynamics, or
multicellular coordination remain visible until discharged by theorem.
Rephrasing does not discharge them.
\end{standingrule}

\begin{standingrule}[\status{D}]\label{sr:bio-no-fitness-primitive}
\textup{(No Fitness Primitive.)}
No fitness landscape, adaptive value, or selection coefficient may be
treated as a primitive. Fitness, to the extent it appears at all,
must emerge as a quotient-visible differential replication rate
licensed by certified replication and hereditary transport structure.
\end{standingrule}

\begin{standingrule}[\status{D}]\label{sr:bio-continuum-gate}
\textup{(Continuum Gate.)}
Continuous-time population dynamics, thermodynamic entropy quantities,
information-theoretic rates, and chemical-kinetics differential
equations may appear only after a branch-local Book~VI bridge theorem
certifies that they faithfully encode certified discrete biological
dynamics within a declared continuum interface regime.
\end{standingrule}

\subsection{Declared Target Regime}\label{sec:bio-target-regime}

\begin{definition}[\status{D}]\label{def:bio-target-regime}
\textup{(Biology Target Regime.)}
A \emph{biology target regime} is a declared branch regime
$\mathcal{R}_{\mathrm{Bio}}$ consisting of:
\begin{enumerate}[label=\textup{(\arabic*)}]
 \item a certified observable class $\mathcal{C}_{\mathrm{obs}}$ of
 physical configurations --- the objects on which replication,
 metabolic, and hereditary probe tests are performed;
 \item a finite admissible probe discipline $\mathcal{P}$, each element
 of which is a declared admissible process detecting a
 quotient-visible distinction among configurations;
 \item a declared resource-transformation window $\mathcal{W}_{\mathrm{met}}$
 specifying the range of metabolic operations within which
 resource-transformation claims have force; and
 \item a declared hereditary comparison discipline governing when two
 descendant configurations are equivalent for the purpose of
 hereditary claims.
\end{enumerate}
The role of $\mathcal{R}_{\mathrm{Bio}}$ is to state explicitly where
biological claims are supposed to have force.
\end{definition}

\begin{remark}[\status{R}]\label{rem:bio-config-note}
A \emph{physical configuration} in this branch is not yet a cell,
bacterium, or organism. At this stage it is only a certified
observable arrangement of matter-like objects on which admissible
probe tests can be performed. Whether any such configuration
qualifies as a replicating entity is determined by the replication
kernel, not by prior biological classification.
\end{remark}

\subsection{Biology Observable Family}\label{sec:bio-obs-family}

\begin{definition}[\status{D}]\label{def:bio-obs-family}
\textup{(Biology Observable Family.)}
A \emph{biology observable family} on $\mathcal{R}_{\mathrm{Bio}}$
is a triple
\[
 \mathcal{O}_{\mathrm{Bio}} \;=\;
 \bigl(\Phi_{\mathrm{Bio}},\;
 \mathcal{M}_{\mathrm{Bio}},\;
 \mathcal{H}_{\mathrm{Bio}}\bigr)
\]
where:
\begin{enumerate}[label=\textup{(\arabic*)}]
 \item $\Phi_{\mathrm{Bio}}$ is the \emph{replication kernel}: the
 quotient-visible self-copy data induced on $\mathcal{C}_{\mathrm{obs}}$
 by admissible replication-probe tests; it records which
 configurations produce certified copies of their own equivalence
 class;
 \item $\mathcal{M}_{\mathrm{Bio}}$ is the \emph{metabolic ledger}:
 certified records of resource-transformation operations --- the
 admissible processes by which a configuration acquires, converts,
 and dissipates resources within the declared metabolic window,
 with full transport bookkeeping; and
 \item $\mathcal{H}_{\mathrm{Bio}}$ is the \emph{heredity-descent
 object}: a branch-visible certified transport map governing how
 observable properties are transmitted from a configuration to its
 certified replication products across declared descent steps.
\end{enumerate}
\end{definition}

\begin{definition}[\status{D}]\label{def:bio-replication-kernel}
\textup{(Replication Kernel.)}
The \emph{replication kernel} $\Phi_{\mathrm{Bio}}$ is the collection
of ordered pairs $(X, X')$ of configurations together with the
certified set of admissible replication-probe tests witnessing that
$X'$ is a copy of $X$ within the declared regime. Two configurations
$X$ and $X'$ satisfy the \emph{replication relation}
$X \xrightarrow{\mathrm{rep}} X'$ if and only if some admissible
probe certifies that $X'$ shares the quotient-visible equivalence
class of $X$ and was produced by an admissible replication process.
\end{definition}

\begin{remark}[\status{R}]\label{rem:bio-rep-remark}
The replication relation is not defined by molecular mechanism,
DNA sequence, or template-copying chemistry. It is defined by
admissible probe certification that the product configuration falls
in the same observable equivalence class as the source configuration.
If, under a given regime, molecular copying turns out to be the
minimal admissible witnessing process, that is a derived theorem, not
a branch assumption.
\end{remark}

\begin{definition}[\status{D}]\label{def:bio-fidelity-margin}
\textup{(Replication Fidelity Margin.)}
The \emph{replication fidelity margin}
\[
 F_{\mathrm{Bio}}(X)
\]
is a branch-specific nonnegative quantity measuring the observable
separation between the equivalence class of a configuration $X$ and
the equivalence classes of its certified replication products. A
positive fidelity margin means the copies are observationally
indistinguishable from the parent within the declared regime. This
is the biological analogue of the spectral margin of the SPEC branch
and the contrast margin of the LING branch.
\end{definition}

\begin{remark}[\status{R}]\label{rem:bio-fidelity-gap-analogy}
The replication fidelity margin is the biological instance of the
canonical gap structure: YM mass gap (spectral separation), SPEC
spectral margin (resonance isolation), LING contrast margin (phonemic
separation), BIO fidelity margin (copy fidelity). In each case the
structural claim is a positive separation between inequivalent
observable classes on the declared regime.
\end{remark}

\begin{definition}[\status{D}]\label{def:bio-metabolic-ledger}
\textup{(Metabolic Ledger.)}
The \emph{metabolic ledger} $\mathcal{M}_{\mathrm{Bio}}$ is a
certified bookkeeping object assigning to each admissible
resource-transformation operation a record of:
\begin{enumerate}[label=\textup{(\arabic*)}]
 \item the resource equivalence classes consumed (input);
 \item the resource equivalence classes produced (output); and
 \item the transport cost and defect contribution incurred, governed
 by the Book~III unified coercive-transport inequality.
\end{enumerate}
The metabolic ledger is additive over declared transformation steps.
\end{definition}

\begin{remark}[\status{R}]\label{rem:bio-metabolism-remark}
The metabolic ledger does not assume thermodynamic equilibrium,
free-energy minimization, or any imported chemical-kinetics framework.
Whether such frameworks are lawfully licensed here is a downstream
Book~VI bridge question. At this stage the metabolic ledger is
a purely discrete accounting object.
\end{remark}

\begin{definition}[\status{D}]\label{def:bio-heredity-descent}
\textup{(Heredity-Descent Object.)}
A \emph{heredity-descent object} $\mathcal{H}_{\mathrm{Bio}}$ is a
branch-visible certified transport map that:
\begin{enumerate}[label=\textup{(\arabic*)}]
 \item is sourced through $\mathbf{U}$;
 \item assigns to each certified replication step
 $X \xrightarrow{\mathrm{rep}} X'$ a declared
 observable-class transmission record; and
 \item produces only quotient-visible distinctions among descendant
 configurations, with all deviations from perfect fidelity accounted
 for within the declared replication fidelity margin.
\end{enumerate}
$\mathcal{H}_{\mathrm{Bio}}$ is not a genetic code, a chromosome map,
or any other pre-given hereditary formalism. It is a derived certified
transport object, admitted only when the branch-legitimacy conditions
of Book~V are satisfied.
\end{definition}

\begin{definition}[\status{D}]\label{def:bio-boundary-object}
\textup{(Biological Boundary Object.)}
A \emph{biological boundary object}
\[
 \mathcal{B}_{\mathrm{Bio}}(X)
\]
is a branch-visible certified interface structure associated with a
configuration $X$, satisfying the Book~II boundary capacity law:
information capacity at the interface is logarithmically bounded by
the declared interface scale. $\mathcal{B}_{\mathrm{Bio}}(X)$
certifies the observable separation between an interior region and an
exterior region for configuration $X$ without importing any background
geometry.
\end{definition}

\begin{remark}[\status{R}]\label{rem:bio-membrane-remark}
The biological boundary object is the canonical NFC surrogate for a
cell membrane. It is not assumed to be a lipid bilayer or any other
chemically specified structure. The only certified property is that
it satisfies the Book~II boundary capacity law and produces a
quotient-visible interior/exterior distinction on the declared regime.
\end{remark}

\subsection{Endpoint Datum}\label{sec:bio-endpoint}

\begin{definition}[\status{D}]\label{def:bio-endpoint}
\textup{(Biology Endpoint Datum.)}
The \emph{biology endpoint datum} is the branch-visible package
\[
 E_{\mathrm{Bio}} \;=\;
 \bigl(\mathcal{C}_{\mathrm{obs}},\;
 \mathcal{O}_{\mathrm{Bio}},\;
 F_{\mathrm{Bio}},\;
 \mathcal{B}_{\mathrm{Bio}},\;
 \mathcal{W}_{\mathrm{met}}\bigr)
\]
consisting of the certified observable class, the biology observable
family, the replication fidelity margin, the biological boundary
object, and the declared metabolic window.
\end{definition}

\begin{definition}[\status{D}]\label{def:lawful-bio-descendant}
\textup{(Lawful Biology Descendant Claim.)}
A \emph{lawful biology descendant claim} is a claim that the endpoint
datum $E_{\mathrm{Bio}}$ is sourced through $\mathbf{U}$ and supports
a certified replication, metabolism, heredity, or boundary statement
on the declared regime.
\end{definition}

\subsection{Current Branch Posture}\label{sec:bio-posture}

\begin{proposition}[\status{C}]\label{prop:bio-branch-candidate}
\textup{(Biology Branch Candidate.)}
\textup{[Conditional on the biology target regime, observable family,
and endpoint datum all being sourced through $\mathbf{U}$,
branch-visible, and admitting a Book~V legitimacy witness.]}
The biology program is a lawful branch candidate.
\end{proposition}

\begin{proof}
Immediate from the branch-legitimacy doctrine (Book~V) applied to the
declared biology endpoint datum.
\qed
\end{proof}

\begin{remark}[\status{R}]\label{rem:bio-posture-note}
Prop.~\ref{prop:bio-branch-candidate} constitutes the branch
architecturally. It does not yet certify any replication theorem,
metabolic theorem, hereditary transport theorem, or boundary
self-organization theorem.
\end{remark}

%% ---------------------------------------------------------------
\section{Admissible Probes and Replication Signatures}
\label{sec:bio-probes}
%% ---------------------------------------------------------------

\subsection{Admissible Replication Probes}\label{sec:bio-admissible-probes}

\begin{definition}[\status{D}]\label{def:bio-admissible-probe}
\textup{(Admissible Replication Probe.)}
An \emph{admissible replication probe} is a declared admissible
process $P \in \mathcal{P}$ whose effect on the certified observable
class is branch-visible and accountable within the spine's
admissible-test discipline. Paradigm cases include copy-verification
tests, boundary-integrity tests, and resource-consumption tests ---
but none of these is assumed as a primitive; each must be declared
and certified within the NFC framework.
\end{definition}

\begin{standingrule}[\status{D}]\label{sr:bio-no-hidden-probe}
\textup{(No Hidden Probe.)}
No admissible replication probe may introduce hidden molecular labels,
undeclared biochemical representations, or non-certified side channels.
Any probe violating this rule is not admissible for biology branch
claims.
\end{standingrule}

\begin{definition}[\status{D}]\label{def:bio-replication-episode}
\textup{(Replication Episode.)}
A \emph{replication episode} beginning from configuration $X$ is a
certified finite sequence of admissible processes --- resource
acquisition, transformation, copy-assembly, boundary-formation ---
producing a configuration $X'$, together with the declared transport
cost record inherited from the metabolic ledger.
\end{definition}

\begin{remark}[\status{R}]\label{rem:bio-episode-note}
A replication episode is the biology analogue of a probe episode
(SPEC branch) or contrast probe episode (LING branch). It is not
a description of a biochemical reaction sequence; it is a certified
admissible process record within the NFC framework.
\end{remark}

\subsection{Replication Signatures and Equivalence}\label{sec:bio-signatures}

\begin{definition}[\status{D}]\label{def:bio-rep-signature}
\textup{(Replication Signature.)}
A \emph{replication signature}
\[
 \mathsf{Rep}(X; t)
\]
is a branch-visible assignment from a configuration $X$ and declared
replication step $t$ to a quotient-visible copy record, together with
its transport and fidelity-defect bookkeeping.
\end{definition}

\begin{standingrule}[\status{D}]\label{sr:bio-signature-invariance}
\textup{(Signature Invariance.)}
A replication signature has theorem-bearing force only to the extent
that it is invariant under presentation-level differences already
eliminated by the observational quotient. Two configurations that
are observationally equivalent cannot bear different replication
signatures.
\end{standingrule}

\begin{definition}[\status{D}]\label{def:bio-rep-equivalence}
\textup{(Replication-Signature Equivalence.)}
Two replication signatures are \emph{replication-equivalent} if they
determine the same quotient-visible copy record under the declared
hereditary comparison discipline of $\mathcal{R}_{\mathrm{Bio}}$.
\end{definition}

\subsection{One-Step Hereditary Ledger}\label{sec:bio-hereditary-ledger}

\begin{definition}[\status{D}]\label{def:bio-one-step-heredity}
\textup{(One-Step Hereditary Ledger Entry.)}
Given a certified replication episode from $X$ to $X'$, the
\emph{one-step hereditary ledger entry}
\[
 \Delta_{\mathrm{her}}\!\bigl(X \xrightarrow{\mathrm{rep}} X'\bigr)
\]
records the certified observable-class transmission from $X$ to $X'$,
together with any declared fidelity deviation, transport cost, or
visible defect contribution inherited from the metabolic ledger.
\end{definition}

\begin{remark}[\status{R}]\label{rem:bio-heredity-note}
The one-step hereditary ledger entry is the biology analogue of the
one-step compositional ledger entry (LING branch) and the one-step
spectroscopy transition ledger (SPEC branch). In each case the core
structure is: take a certified operation, record its effect on the
observable equivalence class, account for the transport cost.
\end{remark}

%% ---------------------------------------------------------------
\section{Core Theorem Shells and Named Frontiers}
\label{sec:bio-theorems}
%% ---------------------------------------------------------------

The following theorem shells articulate the main claims of the
biology branch. Each carries the status it would bear \emph{if}
its stated open obligations were discharged. The open obligations
are named explicitly as branch frontiers.

\subsection{Replication Certification}\label{sec:bio-REP}

\begin{openobligation}[\status{C}]\label{ob:bio-REP}
\textup{(O-BIO.REP: Replication Certification.)}
Show that the replication kernel $\Phi_{\mathrm{Bio}}$ is sourced
through $\mathbf{U}$ and constitutes a lawful branch observable within
the declared regime, with all probe declarations inherited from the
Book~I admissible-test discipline.
\end{openobligation}

\begin{theorem}[\status{C}]\label{thm:bio-rep-certification}
\textup{(Replication Certification Theorem.)}
\textup{[Conditional on O-BIO.REP: Replication Certification.]}
The replication kernel $\Phi_{\mathrm{Bio}}$ is a lawful branch
observable descended from $\mathbf{U}$. The certified replication
relation $X \xrightarrow{\mathrm{rep}} X'$ is quotient-visible and
admissibly testable on the declared regime.
\end{theorem}

\begin{proof}
By the branch-legitimacy doctrine (Book~V) applied to the declared
replication kernel, together with the Book~I observational quotient
discipline. Full proof is conditional on O-BIO.REP.
\qed
\end{proof}

\subsection{Metabolic Ledger Theorem}\label{sec:bio-MET}

\begin{openobligation}[\status{C}]\label{ob:bio-MET}
\textup{(O-BIO.MET: Metabolic Bookkeeping.)}
Show that the metabolic ledger $\mathcal{M}_{\mathrm{Bio}}$ is
additive, sourced through the Book~III transport accounting, and that
all resource-transformation operations preserve transport bookkeeping
without hidden defect.
\end{openobligation}

\begin{theorem}[\status{C}]\label{thm:bio-metabolic-ledger}
\textup{(Metabolic Ledger Theorem.)}
\textup{[Conditional on O-BIO.MET: Metabolic Bookkeeping.]}
The metabolic ledger $\mathcal{M}_{\mathrm{Bio}}$ is additive over
declared resource-transformation steps: transport costs distribute
across components without hidden channel.
\end{theorem}

\begin{proof}
By the Book~III unified coercive-transport inequality applied to the
resource-transformation operations, together with the no-hidden-probe
standing rule (Sr.~\ref{sr:bio-no-hidden-probe}). Full proof is
conditional on O-BIO.MET.
\qed
\end{proof}

\subsection{Replication Fidelity Margin Theorem}\label{sec:bio-FID}

\begin{openobligation}[\status{C}]\label{ob:bio-FID}
\textup{(O-BIO.FID: Replication Fidelity Margin.)}
Show that the replication fidelity margin $F_{\mathrm{Bio}}$ is
branch-computably positive for certified replicating configurations
within the declared regime, and that no fidelity collapse occurs
without a declared admissible error event.
\end{openobligation}

\begin{theorem}[\status{C}]\label{thm:bio-fidelity-margin}
\textup{(Replication Fidelity Margin Theorem.)}
\textup{[Conditional on O-BIO.FID: Replication Fidelity Margin.]}
For any certified replicating configuration $X$, $F_{\mathrm{Bio}}(X)
> 0$: the replication products of $X$ are observationally
indistinguishable from $X$ within the declared fidelity tolerance.
The certified replication inventory is discretely separated from
non-replicating configurations on the declared regime.
\end{theorem}

\begin{proof}
By the observational quotient discipline (Book~I) and the declared
admissible-test finiteness condition on $\mathcal{R}_{\mathrm{Bio}}$.
Positivity of the fidelity margin reflects the separating power of the
admissible replication probe discipline. Full proof is conditional on
O-BIO.FID.
\qed
\end{proof}

\subsection{No-Hidden-Channel Theorem}\label{sec:bio-NC}

\begin{openobligation}[\status{C}]\label{ob:bio-NC}
\textup{(O-BIO.NC: no-hidden Hereditary Channel.)}
Show that no hereditary transport pathway exists outside the declared
one-step hereditary ledger: any discrepancy between two descendant
configurations sharing the same certified parent would require an
undeclared hereditary channel, ruled out by the exhaustive
source-image condition.
\end{openobligation}

\begin{theorem}[\status{C}]\label{thm:bio-no-hidden-channel}
\textup{(No-Hidden-Hereditary-Channel Theorem.)}
\textup{[Conditional on O-BIO.NC: no-hidden Hereditary Channel.]}
Two descendant configurations sharing the same certified parent and
the same admissible replication process have the same replication
signature, up to the declared fidelity tolerance. No undeclared
hereditary pathway can produce a branch-visible distinction invisible
to the declared hereditary ledger.
\end{theorem}

\begin{proof}
Any discrepancy would require an undeclared dissipative or
combinatorial hereditary channel. By the exhaustive source-image
condition (\texttt{def:exhaustive-source-image}, Book~IV) applied to
$\mathbf{U}$ and the no-hidden-gene standing rule
(Sr.~\ref{sr:bio-no-hidden-gene}), no such channel is admitted.
Full proof is conditional on O-BIO.NC.
\qed
\end{proof}

\subsection{Boundary Self-Organization Theorem}\label{sec:bio-BND}

\begin{openobligation}[\status{C}]\label{ob:bio-BND}
\textup{(O-BIO.BND: Boundary Self-Organization.)}
Show that the biological boundary object $\mathcal{B}_{\mathrm{Bio}}(X)$
can emerge from the replication kernel and metabolic ledger without
importing a background spatial geometry: that a certified
configuration can develop a quotient-visible interior/exterior
separation satisfying the Book~II boundary capacity law, as a derived
consequence of replication and metabolic structure.
\end{openobligation}

\begin{theorem}[\status{C}]\label{thm:bio-boundary}
\textup{(Boundary Self-Organization Theorem.)}
\textup{[Conditional on O-BIO.BND: Boundary Self-Organization.]}
For a certified replicating configuration $X$ with positive fidelity
margin and additive metabolic ledger, a biological boundary object
$\mathcal{B}_{\mathrm{Bio}}(X)$ exists satisfying the Book~II
boundary capacity law. The interior/exterior separation is
quotient-visible and sourced through $\mathbf{U}$.
\end{theorem}

\begin{proof}
By the Book~II boundary capacity law applied to the declared interface
scale of $X$, together with the replication certification
(Thm.~\ref{thm:bio-rep-certification}) and metabolic ledger
(Thm.~\ref{thm:bio-metabolic-ledger}). Full proof is conditional on
O-BIO.BND.
\qed
\end{proof}

\begin{remark}[\status{R}]\label{rem:bio-bnd-note}
The boundary self-organization theorem is the canonical NFC surrogate
for the emergence of a cell membrane. It is not a theorem about
lipid bilayers; it is a theorem about the existence of a certified
interface structure satisfying the Book~II capacity law. Whether
lipid bilayers are the minimal physical realization of such an
interface in the target regime is a downstream question, not a branch
assumption.
\end{remark}

\subsection{Hereditary Descent Theorem}\label{sec:bio-HER}

\begin{openobligation}[\status{C}]\label{ob:bio-HER}
\textup{(O-BIO.HER: Hereditary Descent.)}
Show that the heredity-descent object $\mathcal{H}_{\mathrm{Bio}}$ is
sourced through $\mathbf{U}$, passes the five-condition Book~V
branch-legitimacy test, and transmits observable properties across
certified replication steps without hidden channel and within the
declared fidelity margin.
\end{openobligation}

\begin{theorem}[\status{C}]\label{thm:bio-hereditary-descent}
\textup{(Hereditary Descent Theorem.)}
\textup{[Conditional on O-BIO.HER: Hereditary Descent.]}
The heredity-descent object $\mathcal{H}_{\mathrm{Bio}}$ is a lawful
branch-visible certified transport map. Observable properties are
transmitted from $X$ to $X'$ across certified replication steps
within the declared fidelity margin, with no hidden hereditary
channel.
\end{theorem}

\begin{proof}
By the Book~V branch-legitimacy doctrine applied to
$\mathcal{H}_{\mathrm{Bio}}$, together with the no-hidden-gene and
no-smuggling standing rules. Full proof is conditional on O-BIO.HER.
\qed
\end{proof}

\subsection{Evolutionary Dynamics}\label{sec:bio-EVO}

\begin{openobligation}[\status{C}]\label{ob:bio-EVO}
\textup{(O-BIO.EVO: Evolutionary Dynamics.)}
Show that differential replication rates across equivalence classes
of replicating configurations --- a derived quotient-visible quantity
--- give rise to certified selection dynamics within the declared
regime. Fitness is to emerge as a replication-rate differential, not
as an imported adaptive value. This obligation is expected to require
an iterated hereditary ledger and a Book~VI bridge theorem for
continuous-time population dynamics.
\end{openobligation}

\begin{remark}[\status{R}]\label{rem:bio-evo-posture}
\textbf{Posture on evolution.}
Darwinian evolution in the ordinary biological sense refers to
heritable variation, differential fitness, and cumulative change in
population composition over many generations. Within the NFC
framework, fitness cannot be a primitive; it must emerge as a
quotient-visible differential replication rate. Selection dynamics
at the population level, if they appear at all, must be licensed by
a Book~VI bridge theorem as the asymptotic limit of certified
multi-step hereditary transport. Whether such a bridge theorem can
be proved within the NFC framework is a genuine open question.
\end{remark}

\subsection{Multicellular Organization}\label{sec:bio-MULTI}

\begin{openobligation}[\status{C}]\label{ob:bio-MULTI}
\textup{(O-BIO.MULTI: Multicellular Organization.)}
Show, or certify the impossibility of showing within the NFC
framework at the current level, that certified boundary-enclosed
replicating configurations can coordinate into hierarchically
organized multicellular structures --- that is, that
$\mathcal{B}_{\mathrm{Bio}}$ and $\mathcal{H}_{\mathrm{Bio}}$ can
give rise to certified second-order boundary and hereditary objects.
This obligation is deeply downstream of O-BIO.BND and O-BIO.HER and
is not anticipated at the current architectural stage.
\end{openobligation}

\begin{remark}[\status{R}]\label{rem:bio-multi-posture}
\textbf{Posture on multicellularity.}
Multicellular organization is the deepest downstream obligation in
this branch. No claim of multicellular structure is made here.
Multicellularity, if it appears at all, must emerge from iterated
application of the boundary self-organization and hereditary descent
machinery at a higher organizational level. Whether such iteration
is lawfully licensed within the NFC framework is a named open
question, not a currently certified result.
\end{remark}

\subsection{Biology Endpoint}\label{sec:bio-END}

\begin{openobligation}[\status{C}]\label{ob:bio-END}
\textup{(O-BIO.END: Biology Endpoint --- Conditionally Discharged on Nondegenerate Regime $\mathcal{C}^{\mathrm{nd}}_{\mathrm{Bio}}$.)}
Show that the endpoint datum $E_{\mathrm{Bio}}$ is fully certified:
replication certification, metabolic bookkeeping, fidelity margin,
no-hidden-channel, boundary self-organization, and hereditary descent
all discharged. The endpoint is reached when the biology branch
satisfies the shared-endgame schema of Book~VII.
\end{openobligation}

\begin{theorem}[\status{C}]\label{thm:bio-endpoint}
\textup{(Biology Endpoint Theorem.)}
\textup{[Conditional on O-BIO.REP, O-BIO.MET, O-BIO.FID, O-BIO.NC,
O-BIO.BND, O-BIO.HER.]}
The endpoint datum
\[
 E_{\mathrm{Bio}} \;=\;
 \bigl(\mathcal{C}_{\mathrm{obs}},\;
 \mathcal{O}_{\mathrm{Bio}},\;
 F_{\mathrm{Bio}},\;
 \mathcal{B}_{\mathrm{Bio}},\;
 \mathcal{W}_{\mathrm{met}}\bigr)
\]
is a certified lawful biology endpoint: certified replication kernel,
additive metabolic ledger, positive fidelity margin, no-hidden-hereditary-channel control, boundary self-organization sourced through
$\mathbf{U}$, and hereditary descent transport satisfying Book~V
branch-legitimacy. The replication-founded-heredity endpoint is
reached under the stated conditions.
\end{theorem}

\begin{proof}
Combine Thms.~\ref{thm:bio-rep-certification},
\ref{thm:bio-metabolic-ledger}, \ref{thm:bio-fidelity-margin},
\ref{thm:bio-no-hidden-channel}, \ref{thm:bio-boundary}, and
\ref{thm:bio-hereditary-descent}. Apply the shared-endgame schema
(Book~VII, Prop.~6.11): part~(i) is the source-controlled realized
object given by the replication kernel descended from $\mathbf{U}$;
part~(ii) is the metabolic ledger theorem together with the hereditary
descent theorem; part~(iii) is the no-hidden-channel theorem together
with the fidelity margin. Full proof is conditional on O-BIO.REP
through O-BIO.HER.
\qed
\end{proof}

%% ---------------------------------------------------------------
\section{Closure Ledger}\label{sec:bio-closure}
%% ---------------------------------------------------------------

\subsection{Summary of Named Open Obligations}

The following table records all named open obligations in the biology
branch as of the current canonical state.

\medskip
\noindent
\begin{tabular}{|l|l|p{8.0cm}|}
\hline
\textbf{Label} & \textbf{Status} & \textbf{Description} \\
\hline
\texttt{ob:bio-REP} & [O] & Replication Certification: $\Phi_{\mathrm{Bio}}$ sourced through $\mathbf{U}$ \\
\texttt{ob:bio-MET} & [O] & Metabolic Bookkeeping: additivity and transport accounting \\
\texttt{ob:bio-FID} & [O] & Replication Fidelity Margin: positive separation of replicating class \\
\texttt{ob:bio-NC} & [O] & no-hidden Hereditary Channel: no undeclared hereditary pathway \\
\texttt{ob:bio-BND} & [O] & Boundary Self-Organization: interface structure from Book~II capacity law \\
\texttt{ob:bio-HER} & [O] & Hereditary Descent: $\mathcal{H}_{\mathrm{Bio}}$ passes Book~V test \\
\texttt{ob:bio-EVO} & [O] & Evolutionary Dynamics: selection as differential replication rate (Book~VI bridge likely needed) \\
\texttt{ob:bio-MULTI} & [O] & Multicellular Organization: second-order boundary and hereditary objects (deeply downstream) \\
\texttt{ob:bio-END} & [O] & Endpoint: full certification of $E_{\mathrm{Bio}}$ \\
\hline
\end{tabular}

\medskip

\begin{remark}[\status{R}]\label{rem:bio-discharge-order}
\textbf{Discharge order.}
The natural discharge order is: REP $\to$ MET $\to$ FID $\to$ NC
$\to$ BND $\to$ HER $\to$ END. Evolutionary dynamics (EVO) and
multicellular organization (MULTI) are downstream of END and may
require Book~VI bridge theorems. Crucially, END can be reached
without discharging EVO or MULTI: the
replication-founded-heredity endpoint is the minimal honest claim
of this branch.
\end{remark}

\subsection{Two-Status Reading}

\begin{remark}[\status{R}]\label{rem:bio-two-status}
\textbf{Two-status reading.}
The biology branch admits two natural closure readings:
\begin{enumerate}[label=\textup{(\alph*)}]
 \item \emph{Replication-heredity closure} (obligations REP--END
 discharged, EVO and MULTI open): certified replication kernel,
 additive metabolic ledger, positive fidelity margin, boundary
 self-organization, and hereditary descent --- the minimal honest
 biological endpoint within the NFC framework.
 \item \emph{Evolutionary closure} (REP--EVO discharged): adds
 certified selection dynamics as a quotient-visible differential
 replication rate over multiple descent steps, requiring a Book~VI
 bridge for continuous-time population dynamics.
\end{enumerate}
Posture~(a) is the correct near-term target. Posture~(b) is the
genuine long-term ambition for the branch.
\end{remark}

\subsection{Relation to Existing Branches}

\begin{remark}[\status{R}]\label{rem:bio-branch-relations}
\textbf{Relation to existing branches.}
The biology branch imports structural templates from the NS, SPEC, and
LING branches but does not depend on any of their open obligations
being discharged. In particular:
\begin{enumerate}[label=\textup{(\arabic*)}]
 \item The replication-continuation structure is formally parallel to
 the NS continuation criterion (with replication episodes replacing
 differential-equation time steps). No NS branch obligation is a
 prerequisite.
 \item The metabolic ledger is formally parallel to the SPEC
 transition ledger (with resource-transformation replacing
 probe-induced spectral transitions). No SPEC branch obligation
 is a prerequisite.
 \item The no-hidden-channel structure is formally parallel to the
 LING no-hidden-channel theorem. No LING branch obligation is a
 prerequisite.
 \item The boundary self-organization theorem imports only the
 Book~II boundary capacity law from the spine; it does not depend
 on any branch book result.
 \item The YM, GR, SM, RH, and SCC branches have no direct bearing
 on the biology branch at the replication-heredity level. Cross-branch
 obligations may emerge if matter-specificity or physical-law
 constraints are required for the metabolic ledger, but this is not
 anticipated at the current architectural stage.
\end{enumerate}
\end{remark}

%% ---------------------------------------------------------------
\section{Obligation Discharges}\label{sec:bio-discharges}
%% ---------------------------------------------------------------

This section discharges obligations O-BIO.REP through O-BIO.HER by
appeal to existing certified spine results. Each discharge replaces
the corresponding open obligation with a conditional theorem whose
hypotheses are stated explicitly. O-BIO.BND is named as the
\emph{principal biology frontier}: Book~II bounds the capacity of
any self-organized interface but does not derive that a
replicating-metabolic configuration will produce one --- that
derivation is the central open obligation of this branch.
O-BIO.EVO and O-BIO.MULTI remain open. O-BIO.END is blocked by
O-BIO.BND and becomes available only after the principal frontier
is resolved.

\subsection{O-BIO.REP: Replication Certification}\label{sec:bio-REP-discharge}

\begin{theorem}[\status{C}]\label{thm:bio-REP-discharged}
\textup{(Replication Certification --- O-BIO.REP Conditionally Discharged.)}
\textup{[Conditional on: (1) a declared biology target regime
$\mathcal{R}_{\mathrm{Bio}}$ with admissible replication probe family
$\mathcal{P} \subseteq \mathcal{T}_n$ (Book~I admissible tests) and
certified observable class $\mathcal{C}_{\mathrm{obs}}$ descended from
$\mathbf{U}$ (Book~IV, \texttt{thm:universal-completion}); and (2)
Book~II transfer structure (\texttt{def:transfer-system}) on
$\mathcal{C}_{\mathrm{obs}}$.]}
There exists a certified replication kernel
\[
 \Phi_{\mathrm{Bio}} : X \;\mapsto\;
 \bigl\{X' \in \mathcal{C}_{\mathrm{obs}} :
 \exists P \in \mathcal{P},\; P \text{ certifies }
 X \xrightarrow{\mathrm{rep}} X'\bigr\}
\]
where each $X'$ satisfies $[X']_{\sim_n} = [X]_{\sim_n}$
(Book~I observable equivalence) and is produced by a declared
admissible replication process. The replication relation is
quotient-visible, transport-accountable, and defect-tracked.
\end{theorem}

\begin{proof}
The observable equivalence relation $\sim_n$ (Book~I,
\texttt{def:obs-equiv}) is defined on all admissible-test outcomes.
A replication-probe $P \in \mathcal{P}$ certifies $X \to X'$ as a
replication episode if and only if $P$ witnesses
$[X']_{\sim_n} = [X]_{\sim_n}$ and records the production process
as admissible. Restricting to such probes gives the quotient-visible
replication kernel by the no-smuggling standing rule
(Sr.~\ref{sr:bio-no-smuggling}). Transport accounting is provided
by Book~III persistence and transfer structure. The boundary-capacity
bookkeeping from Book~II bounds the information capacity of any
replication interface within $\mathcal{R}_{\mathrm{Bio}}$. No
imported molecular mechanism or genetic code is required.
\qed
\end{proof}

\begin{remark}[\status{R}]\label{rem:bio-REP-scope}
The replication kernel records the structural minimum: which
configurations produce certified same-class copies under admissible
probes. The richer structure --- metabolic accounting, fidelity
margin, hereditary transport --- builds on this foundation in
subsequent discharges.
\end{remark}

\subsection{O-BIO.MET: Metabolic Bookkeeping}\label{sec:bio-MET-discharge}

\begin{theorem}[\status{C}]\label{thm:bio-MET-discharged}
\textup{(Metabolic Bookkeeping --- O-BIO.MET Conditionally Discharged.)}
\textup{[Conditional on: O-BIO.REP discharged
(Thm.~\ref{thm:bio-REP-discharged}); Book~III unified
coercive-transport inequality (\texttt{thm:unified-coercive}, [U])
and ledger additivity (\texttt{lem:transport-equiv-null}).]}
The metabolic ledger $\mathcal{M}_{\mathrm{Bio}}$ is additive over
declared resource-transformation steps. For each admissible
resource-transformation operation, the transport cost is bounded by
the Book~III coercive-transport inequality applied to the
transformation, and costs distribute additively across input and
output resource equivalence classes without hidden channel.
\end{theorem}

\begin{proof}
Each resource-transformation operation is an admissible process on
$\mathcal{C}_{\mathrm{obs}}$, and its input and output resource
classes are quotient-visible by O-BIO.REP. The Book~III unified
coercive-transport inequality bounds the transport cost of each
operation. By \texttt{lem:transport-equiv-null}, ledger entries
are invariant under admissible relabelling of resource classes.
Additivity over multi-step transformations follows by induction: each
step contributes its bounded transport cost to the running ledger
total. The no-hidden-probe standing rule (Sr.~\ref{sr:bio-no-hidden-probe})
ensures no undeclared energy or matter channel enters.
\qed
\end{proof}

\subsection{O-BIO.FID: Replication Fidelity Margin}\label{sec:bio-FID-discharge}

\begin{theorem}[\status{C}]\label{thm:bio-FID-discharged}
\textup{(Replication Fidelity Margin --- O-BIO.FID Conditionally Discharged.)}
\textup{[Conditional on: O-BIO.REP discharged
(Thm.~\ref{thm:bio-REP-discharged}); the admissible replication probe
family $\mathcal{P}$ being non-trivially fidelity-certifying on the
declared regime, i.e., every certified replication episode
$X \xrightarrow{\mathrm{rep}} X'$ is witnessed by a probe confirming
$[X']_{\sim_n} = [X]_{\sim_n}$.]}
For every certified replicating configuration $X$,
$F_{\mathrm{Bio}}(X) > 0$: the replication products of $X$ are
observationally indistinguishable from $X$ within the declared fidelity
tolerance. The certified replication inventory is discretely separated
from non-replicating configurations on the declared regime.
\end{theorem}

\begin{proof}
By O-BIO.REP, a replication episode $X \xrightarrow{\mathrm{rep}} X'$
is certified if and only if an admissible probe witnesses
$[X']_{\sim_n} = [X]_{\sim_n}$. The fidelity margin
$F_{\mathrm{Bio}}(X)$ is the quotient-visible measure of the
class-agreement between parent and products
(Def.~\ref{def:bio-fidelity-margin}). Under the fidelity-certifying
condition, every certified replication product falls in the same
equivalence class as its parent; hence $F_{\mathrm{Bio}}(X)$ is
positive by the separating power of the admissible probe family
(same argument as O-LING.CM: distinct classes are separated by
definition). The discrete separation from non-replicating
configurations follows because non-replicating configurations, by
definition, have no admissible probe certifying a same-class copy.
\qed
\end{proof}

\subsection{O-BIO.NC: no-hidden Hereditary Channel}\label{sec:bio-NC-discharge}

\begin{theorem}[\status{C}]\label{thm:bio-NC-discharged}
\textup{(no-hidden Hereditary Channel --- O-BIO.NC Conditionally Discharged.)}
\textup{[Conditional on: O-BIO.MET discharged
(Thm.~\ref{thm:bio-MET-discharged}); two descendant configurations
sharing the same certified parent and the same admissible replication
process; and the universal completion theorem
(\texttt{thm:universal-completion}, Book~IV, [U]) applied via the
exhaustive source-image condition
(\texttt{def:exhaustive-source-image}, Book~IV).]}
Two descendant configurations sharing the same certified parent and
admissible replication process have the same replication signature,
up to the declared fidelity tolerance. No undeclared hereditary
pathway can produce a branch-visible distinction invisible to the
declared hereditary ledger.
\end{theorem}

\begin{proof}
Let $X'$ and $X''$ both be certified replication products of $X$
under the same admissible process. By O-BIO.REP, both satisfy
$[X']_{\sim_n} = [X'']_{\sim_n} = [X]_{\sim_n}$. By O-BIO.MET,
the metabolic ledger entries for both production episodes are
determined by the same input resource classes and assembly operations.
Suppose for contradiction that some probe $P \in \mathcal{P}$
distinguishes $X'$ from $X''$. This witnesses a difference in their
observable equivalence classes not attributable to any declared
parent-level or process-level difference --- an undeclared hereditary
channel. By the exhaustive source-image condition
(\texttt{def:exhaustive-source-image}, Book~IV) applied to
$\mathbf{U}$, no such channel exists. Contradiction; hence
$[X']_{\sim_n} = [X'']_{\sim_n}$ and their replication signatures
agree.
\qed
\end{proof}

\subsection{O-BIO.HER: Hereditary Descent}\label{sec:bio-HER-discharge}

\begin{theorem}[\status{C}]\label{thm:bio-HER-discharged}
\textup{(Hereditary Descent --- O-BIO.HER Conditionally Discharged.)}
\textup{[Conditional on: O-BIO.REP (Thm.~\ref{thm:bio-REP-discharged}),
O-BIO.FID (Thm.~\ref{thm:bio-FID-discharged}), and O-BIO.NC
(Thm.~\ref{thm:bio-NC-discharged}) all discharged; the declared
replication probe family being closed under iterated application in
the sense that replication products are themselves certified observable
configurations subject to the same probe discipline; and Book~V
branch-legitimacy doctrine (\texttt{thm:UBLT}, [U]).]}
There exists a certified heredity-descent object
$\mathcal{H}_{\mathrm{Bio}}$ sourced through $\mathbf{U}$, identified
as the restriction of the Book~III coercive-transport map to the
replication-visible sector of $\mathcal{C}_{\mathrm{obs}}$ on
$\mathcal{R}_{\mathrm{Bio}}$:
\begin{enumerate}[label=\textup{(\arabic*)}]
 \item $\mathcal{H}_{\mathrm{Bio}}$ assigns to each certified
 replication step $X \xrightarrow{\mathrm{rep}} X'$ a declared
 observable-class transmission record;
 \item all fidelity deviations are accounted for within the declared
 fidelity margin;
 \item outputs are quotient-visible within $\mathcal{R}_{\mathrm{Bio}}$;
 \item $\mathcal{H}_{\mathrm{Bio}}$ passes the Book~V five-condition
 branch-legitimacy test.
\end{enumerate}
\end{theorem}

\begin{proof}
By O-BIO.REP, the replication relation is certified and
quotient-visible. The Book~III coercive-transport governing theorem
constructs a canonical transport map on every certified
transport-system descendant of $\mathbf{U}$; restricting to the
replication-visible sector gives $\mathcal{H}_{\mathrm{Bio}}$. By
O-BIO.FID, each replication product satisfies
$[X']_{\sim_n} = [X]_{\sim_n}$ with positive fidelity margin; hence
the transmission record is well-defined. By O-BIO.NC, the
hereditary channel is unique (no-hidden alternatives); hence the
transport map is path-independent on the replication-visible sector.
The Book~V five-condition test is satisfied:
\begin{enumerate}[label=\textup{(\roman*)}]
 \item explicit branch candidate ($\mathcal{H}_{\mathrm{Bio}}$ declared);
 \item observable family visibility (replication-class outputs are
 quotient-visible);
 \item source-descended comparison (inherits from $\mathbf{U}$ via
 Book~III);
 \item no-hidden supplement (Sr.~\ref{sr:bio-no-hidden-gene});
 \item no path multiplication (O-BIO.NC ensures uniqueness of
 observable-class output per replication episode).
\end{enumerate}
\qed
\end{proof}

\begin{remark}[\status{R}]\label{rem:bio-partial-endpoint}
\textbf{Partial endpoint status.}
Under the hypotheses of
Thms.~\ref{thm:bio-REP-discharged}--\ref{thm:bio-HER-discharged},
the biology branch carries: a certified replication kernel; additive
metabolic ledger; positive replication fidelity margin; no-hidden-hereditary-channel control; and a hereditary descent object passing
Book~V legitimacy. This constitutes the \emph{replication-heredity
partial endpoint} --- the maximum honest closure achievable without
resolving the principal frontier (O-BIO.BND). The full endpoint
datum $E_{\mathrm{Bio}}$ requires O-BIO.BND; accordingly O-BIO.END
remains blocked.
\end{remark}

\subsection{O-BIO.BND: Principal Biology Frontier (Remains Open)}
\label{sec:bio-BND-frontier}

\begin{openobligation}[\status{O}]\label{ob:bio-BND-open}
\textup{(O-BIO.BND: Boundary Self-Organization --- Principal Biology
Frontier.)}
The boundary self-organization obligation is the central named
frontier of the biology branch. Book~II establishes a logarithmic
boundary capacity law bounding the information capacity of any
interface, but does not derive that a replicating-metabolic
configuration will spontaneously produce a certified
interior/exterior-separating boundary object. The required
derivation must show that: (1) the replication kernel $\Phi_{\mathrm{Bio}}$
and metabolic ledger $\mathcal{M}_{\mathrm{Bio}}$, together, imply the
existence of a quotient-visible interface structure on some class of
certified configurations; (2) that interface satisfies the Book~II
capacity bound without importing background geometry; and (3) the
interior/exterior separation is sourced through $\mathbf{U}$.
This is not dischargeable by appeal to existing spine results alone;
it requires new canonical content specific to the biology branch.
O-BIO.BND is the canonical analogue of the NS Stage-3 frontier
(thm:NS61): the principal open obligation gating the full endpoint.
\end{openobligation}

\subsection{O-BIO.BND: Boundary Self-Organization — Discharge on Nondegenerate Regime}
\label{sec:bio-BND-discharge}

The principal biology frontier is split into two parts:
$\mathrm{O\text{-}BIO.BND} = \mathrm{O\text{-}BIO.BND}_{\mathrm{subclass}}
+ \mathrm{O\text{-}BIO.BND}_{\mathrm{force}}$.
The first is discharged here on a declared nondegenerate endpoint
regime. The second is named as the new minimal frontier.

\begin{definition}[\status{D}]\label{def:bio-rep-maintaining-core}
\textup{(Replication-Maintaining Core.)}
For a certified replicating configuration $X \in \mathcal{C}_{\mathrm{obs}}$,
the \emph{replication-maintaining core} is
\[
 I_X \;:=\; \bigcap\bigl\{S \subseteq X :
 S \text{ preserves the certified replication signature of } X\bigr\}.
\]
This is not a membrane. It is the minimal quotient-visible support
preserving the already-certified replication signature
(Def.~\ref{def:bio-rep-signature}). Membership in $I_X$ is determined
only by contribution to the certified replication observable, not by
raw spatial naming.
\end{definition}

\begin{definition}[\status{D}]\label{def:bio-boundary-forcing-exchange}
\textup{(Boundary-Forcing Exchange.)}
A certified configuration $X$ has \emph{boundary-forcing exchange}
when there exists a metabolic ledger entry $m \in \mathcal{M}_{\mathrm{Bio}}$
such that
\[
 \mathrm{supp}(m) \cap I_X \neq \varnothing
 \quad\text{and}\quad
 \mathrm{supp}(m) \cap (X \setminus I_X) \neq \varnothing.
\]
At least one declared resource-transformation step touches both the
replication-maintaining core and its complement.
\end{definition}

\begin{definition}[\status{D}]\label{def:bio-nondegenerate-regime}
\textup{(Nondegenerate Certified BIO Regime.)}
Let $\mathcal{C}_{\mathrm{Bio}}^{\mathrm{nd}} \subseteq \mathcal{C}_{\mathrm{obs}}$
be the class of certified BIO configurations $X$ satisfying:
\begin{enumerate}[label=\textup{(\arabic*)}]
 \item certified replication under $\Phi_{\mathrm{Bio}}$;
 \item positive fidelity margin $F_{\mathrm{Bio}}(X) > 0$;
 \item additive metabolic ledger $\mathcal{M}_{\mathrm{Bio}}$
 with at least one declared resource-transformation step;
 \item no-hidden hereditary/metabolic channel;
 \item proper replication-maintaining core:
 $\varnothing \neq I_X \subsetneq X$;
 \item nontrivial metabolic dependence: persistence of the
 replication signature depends on at least one metabolic
 ledger entry;
 \item external resource support: at least one replication-relevant
 resource support lies outside $I_X$.
\end{enumerate}
This imports no biology. It only restricts the endpoint to the class
where replication, metabolism, and external resource exchange are all
genuinely present.
\end{definition}

\begin{theorem}[\status{C}]\label{thm:bio-BND-boundary-forcing}
\textup{(Boundary-Forcing Exchange Lemma.)}
\textup{[Conditional on certified replication (O-BIO.REP discharged),
positive fidelity (O-BIO.FID discharged), additive metabolic
bookkeeping (O-BIO.MET discharged), no-hidden-channel discipline
(O-BIO.NC discharged), and boundary-forcing exchange.]}
There exists a quotient-visible finite interface
\[
 \mathcal{B}_{\mathrm{Bio}}(X) \;:=\; \partial I_X
\]
separating the replication-maintaining core $I_X$ from the exchange
complement $X \setminus I_X$. This interface is sourced through the
declared replication and metabolic observables, hence through
$\mathbf{U}$, and falls under Book~II boundary-capacity control.
\end{theorem}

\begin{proof}
By O-BIO.REP, the replication signature of $X$ is quotient-visible,
determined by admissible replication probes, not by molecular
chemistry. Because $I_X$ is the intersection of all supports
preserving that certified signature, membership in $I_X$ is itself
quotient-visible: it is determined by contribution to the certified
replication observable.

By O-BIO.MET, the metabolic ledger is additive over declared
resource-transformation steps with no-hidden resource channel.
By boundary-forcing exchange, some $m \in \mathcal{M}_{\mathrm{Bio}}$
meets both $I_X$ and $X \setminus I_X$. The partition
$(I_X,\, X \setminus I_X)$ is therefore distinguished by a certified
observable fact: a declared metabolic transformation crosses the
partition while the replication signature is maintained by $I_X$.

If the crossing were not quotient-visible, $m$ would either be absent
from the declared metabolic ledger or would require a hidden exchange
channel, contradicting O-BIO.MET and O-BIO.NC. Hence
$\partial I_X$ is a certified interface. Since it is a finite
relation-neighborhood defined from certified BIO observables, the
Book~II boundary capacity law applies once it is recognized as an
interface. The BIO branch already says Book~II bounds an interface
once one exists; the missing part was deriving the interface from
replication-metabolism structure. That derivation is supplied by
boundary-forcing exchange.
\qed
\end{proof}

\begin{theorem}[\status{C}]\label{thm:bio-BND-closed-core-alt}
\textup{(Closed-Core Alternative.)}
\textup{[Conditional on certified replication, positive fidelity,
additive metabolic ledger, no-hidden channel.]}
If no metabolic ledger entry crosses the replication-maintaining core
boundary $\partial I_X$, then exactly one of the following holds:
\begin{enumerate}[label=\textup{(\arabic*)}]
 \item \emph{Internal-only metabolism}: every
 replication-relevant transformation is supported inside $I_X$;
 \item \emph{External-only metabolism}: every transformation is
 supported outside $I_X$ and irrelevant to replication
 maintenance;
 \item \emph{Trivial whole-support core}: $I_X = X$;
 \item \emph{Hidden-channel violation}: replication persistence
 depends on resource transformation not recorded in
 $\mathcal{M}_{\mathrm{Bio}}$.
\end{enumerate}
Since option~4 is forbidden by no-hidden-channel discipline (O-BIO.NC),
the only no-crossing cases are boundary-degenerate configurations.
\end{theorem}

\begin{proof}
Each metabolic ledger entry has finite support inside the declared
resource-transformation window $\mathcal{W}_{\mathrm{met}}$. If no
entry crosses $\partial I_X$, each $m \in \mathcal{M}_{\mathrm{Bio}}$
has support entirely in $I_X$, entirely in $X \setminus I_X$, or
outside the replication-relevant support. This yields exactly the
four cases. Option~4 requires an undeclared channel, excluded by
O-BIO.NC and the exhaustive source-image condition
(\texttt{def:exhaustive-source-image}, Book~IV) applied to $\mathbf{U}$.
\qed
\end{proof}

\begin{corollary}[\status{C}]\label{cor:bio-BND-nondegenerate-necessity}
\textup{(Nondegenerate Boundary Necessity.)}
\textup{[Conditional on the no-crossing alternatives of
Thm.~\ref{thm:bio-BND-closed-core-alt} and the nondegeneracy
conditions of Def.~\ref{def:bio-nondegenerate-regime}.]}
Every $X \in \mathcal{C}_{\mathrm{Bio}}^{\mathrm{nd}}$ satisfies
boundary-forcing exchange. Hence
$\mathcal{B}_{\mathrm{Bio}}(X) = \partial I_X$ exists as a
quotient-visible finite interface sourced through $\mathbf{U}$,
and Book~II boundary-capacity control applies.
\end{corollary}

\begin{proof}
Assume for contradiction that no metabolic entry crosses
$\partial I_X$. By Thm.~\ref{thm:bio-BND-closed-core-alt}, the
configuration falls into one of four cases. The nondegeneracy
conditions exclude all four:
$I_X \subsetneq X$ excludes whole-support;
external resource support excludes internal-only;
nontrivial metabolic dependence excludes external-only irrelevance;
O-BIO.NC excludes hidden-channel violation. Contradiction; hence
boundary-forcing exchange holds and
Thm.~\ref{thm:bio-BND-boundary-forcing} applies.
\qed
\end{proof}

\begin{theorem}[\status{C}]\label{thm:bio-BND-nd}
\textup{(Nondegenerate BIO Boundary Theorem.)}
\textup{[Conditional on O-BIO.REP through O-BIO.HER all discharged
and the nondegeneracy conditions of Def.~\ref{def:bio-nondegenerate-regime}.]}
For every $X \in \mathcal{C}_{\mathrm{Bio}}^{\mathrm{nd}}$, there
exists a quotient-visible finite boundary object
$\mathcal{B}_{\mathrm{Bio}}(X) = \partial I_X$ separating the
replication-maintaining core $I_X$ from its exchange complement.
O-BIO.BND is discharged on $\mathcal{C}_{\mathrm{Bio}}^{\mathrm{nd}}$.
\end{theorem}

\begin{proof}
By Cor.~\ref{cor:bio-BND-nondegenerate-necessity}, every
$X \in \mathcal{C}_{\mathrm{Bio}}^{\mathrm{nd}}$ satisfies
boundary-forcing exchange. Apply
Thm.~\ref{thm:bio-BND-boundary-forcing}: $\partial I_X$ is a
certified quotient-visible interface sourced through $\mathbf{U}$
with Book~II capacity control. This is precisely the content of
O-BIO.BND restricted to the nondegenerate regime.
\qed
\end{proof}

\begin{remark}[\status{R}]\label{rem:bio-BND-force-frontier}
\textbf{Remaining frontier: O-BIO.BND-Force.}
The discharge of O-BIO.BND on $\mathcal{C}_{\mathrm{Bio}}^{\mathrm{nd}}$
does not prove that every abstract certified replicator is nondegenerate.
The remaining minimal frontier is:
\[
 \mathrm{O\text{-}BIO.BND}_{\mathrm{force}}:
 \quad\text{prove that every nontrivial persistent
 replication-metabolism cycle has } I_X \subsetneq X
 \text{ and external resource support.}
\]
In plain terms: can a nontrivial persistent certified
replication-metabolism cycle avoid all exchange across $\partial I_X$?
This is the new minimal biology frontier.
\end{remark}

\begin{theorem}[\status{C}]\label{thm:bio-END-nd}
\textup{(Nondegenerate Biology Endpoint.)}
\textup{[Conditional on O-BIO.REP through O-BIO.HER discharged and
$X \in \mathcal{C}_{\mathrm{Bio}}^{\mathrm{nd}}$.]}
The endpoint datum
\[
 E_{\mathrm{Bio}}^{\mathrm{nd}} \;=\;
 \bigl(\mathcal{C}_{\mathrm{Bio}}^{\mathrm{nd}},\;
 \mathcal{O}_{\mathrm{Bio}},\;
 F_{\mathrm{Bio}},\;
 \mathcal{B}_{\mathrm{Bio}},\;
 \mathcal{W}_{\mathrm{met}}\bigr)
\]
is a certified lawful biology endpoint at the replication-heredity
level on the declared nondegenerate regime.
\end{theorem}

\begin{proof}
The five discharged obligations (REP, MET, FID, NC, HER) supply:
certified replication kernel, additive metabolic ledger, positive
fidelity margin, no-hidden-channel control, and hereditary descent
sourced through $\mathbf{U}$. By Thm.~\ref{thm:bio-BND-nd},
O-BIO.BND is discharged on $\mathcal{C}_{\mathrm{Bio}}^{\mathrm{nd}}$.
Apply the Book~VII shared-endgame schema (Prop.~6.11) to
$E_{\mathrm{Bio}}^{\mathrm{nd}}$: part~(i) is the replication kernel
from $\mathbf{U}$; part~(ii) is the metabolic ledger and hereditary
descent; part~(iii) is the no-hidden-channel theorem with the fidelity
margin. The boundary object from $\mathbf{U}$ and Book~II capacity
close the package.
\qed
\end{proof}

%% ---------------------------------------------------------------
\section{Discrete Evolutionary Dynamics}\label{sec:bio-EVO-discrete}
%% ---------------------------------------------------------------

\begin{definition}[\status{D}]\label{def:bio-iterated-hereditary-ledger}
\textup{(Iterated Hereditary Ledger.)}
On $\mathcal{C}_{\mathrm{Bio}}^{\mathrm{nd}}$, the
\emph{$n$-step iterated hereditary ledger}
\[
 \mathcal{H}_n(X)
\]
is the finite certified record of all replication-descendant classes
reachable from $X$ in $n$ declared replication steps, together with
their certified transport and defect bookkeeping.
\end{definition}

\begin{definition}[\status{D}]\label{def:bio-replication-rate}
\textup{(Finite-Stage Replication Rate.)}
For a hereditary equivalence class $[X]$, define its $n$-step
certified replication count
\[
 N_n([X]) \;:=\;
 \#\bigl\{X' \in \mathcal{C}_{\mathrm{Bio}}^{\mathrm{nd}} :
 X \xrightarrow{n\text{-rep}} X'
 \text{ by certified hereditary descent}\bigr\}
\]
and the \emph{finite-stage replication rate}
$r_n([X]) := \tfrac{1}{n}\log N_n([X])$
when $N_n([X]) > 0$.
This is not a primitive fitness; it is a derived quotient-visible
count over certified hereditary descendants.
\end{definition}

\begin{definition}[\status{D}]\label{def:bio-differential-replication-advantage}
\textup{(Differential Replication Advantage.)}
For two hereditary classes $[X]$ and $[Y]$, the
\emph{finite-stage selection advantage} of $[X]$ over $[Y]$ at
stage $n$ is
\[
 \Delta r_n([X],[Y]) \;:=\; r_n([X]) - r_n([Y]).
\]
A class $[X]$ has finite-stage selection advantage if $\Delta r_n > 0$.
\end{definition}

\begin{theorem}[\status{C}]\label{thm:bio-EVO-discrete-selection}
\textup{(Discrete Selection Visibility.)}
\textup{[Conditional on O-BIO.REP, O-BIO.HER discharged and
$\mathcal{C}_{\mathrm{Bio}}^{\mathrm{nd}}$ declared.]}
Finite-stage differential replication advantage
$\Delta r_n([X],[Y])$ is quotient-visible and source-descended.
\end{theorem}

\begin{proof}
The replication kernel $\Phi_{\mathrm{Bio}}$ records certified
same-class copy events; the hereditary object $\mathcal{H}_{\mathrm{Bio}}$
transmits observable properties with no-hidden channel. Therefore
$N_n([X])$ is computed only from certified replication and hereditary
data: the descendant count is ledger-visible. Since $[X]$ and $[Y]$
are observational equivalence classes and their descendant records are
quotient-visible, $\Delta r_n([X],[Y])$ is quotient-visible. Since all
components descend from the certified BIO endpoint data through
$\mathbf{U}$, the replication-rate differential is source-descended.
\qed
\end{proof}

\begin{theorem}[\status{C}]\label{thm:bio-EVO-fitness-as-rate}
\textup{(Fitness-as-Rate Theorem.)}
\textup{[Conditional on Thm.~\ref{thm:bio-EVO-discrete-selection}.]}
Define the \emph{finite-stage fitness surrogate}
$\mathrm{Fit}_n([X]) := r_n([X])$.
Then $\mathrm{Fit}_n$ is a lawful finite-stage fitness surrogate: it is
quotient-visible, hereditary-ledger-derived, and not imported as an
adaptive value. Standing Rule~\ref{sr:bio-no-fitness-primitive}
is satisfied.
\end{theorem}

\begin{proof}
Fitness is defined as quotient-visible replication rate, not as a
primitive adaptive quantity. Thm.~\ref{thm:bio-EVO-discrete-selection}
provides this. Sr.~\ref{sr:bio-no-fitness-primitive} forbids primitive
fitness and requires it to emerge as differential replication rate;
$\mathrm{Fit}_n$ satisfies exactly that criterion.
\qed
\end{proof}

\begin{theorem}[\status{C}]\label{thm:bio-EVO-discrete-dynamics}
\textup{(Discrete Selection Dynamics.)}
\textup{[Conditional on Thm.~\ref{thm:bio-EVO-fitness-as-rate}
and the iterated hereditary ledger being closed under $n \mapsto n+1$.]}
Finite-stage population composition evolves by certified differential
replication rates. The one-step transition is
\[
 p_{n+1}([X]) \;=\;
 \frac{p_n([X])\,N_1([X])}{\sum_{[Y]} p_n([Y])\,N_1([Y])}
\]
or, more generally, by the declared finite-step hereditary transition
matrix $P_n([X] \to [Y])$ with no fitness landscape or external
population law imported.
\end{theorem}

\begin{proof}
Transition weights are computed from certified descendant counts and
hereditary transport records; normalization is bookkeeping over
quotient-visible classes. No fitness landscape or external
population law is imported; the update rule is the finite ledger
summary of certified replication outcomes.
\qed
\end{proof}

\begin{openobligation}[\status{C}]\label{ob:bio-EVO-pop-bridge}
\textup{(O-BIO.EVO-Bridge: Population Dynamics Book~VI Bridge.)}
\textup{[Conditional on Thm.~\ref{thm:bio-EVO-discrete-dynamics}
and the B2-DomainNoLoss template below.]}
Continuous-time selection dynamics require a Book~VI bridge theorem.
The following theorem reduces the bridge to a single named
convergence condition.
\end{openobligation}

\begin{definition}[\status{D}]\label{def:bio-EVO-domain-loss}
\textup{(Evolutionary Domain Loss.)}
For the discrete hereditary transition family $(P_n)_{n\ge0}$
and a declared continuous-time comparison dynamics $\Phi_t$,
the \emph{evolutionary domain loss} at stage $n$ is
\[
 \varepsilon_n^{\mathrm{evo}} :=
 \sup_{\psi\,\mathrm{admissible}}\;
 \bigl|
 \mathcal{R}_{\mathrm{disc},n}(\psi)
 - \mathcal{R}_{\mathrm{cts},n}(\psi)
 \bigr|,
\]
where $\mathcal{R}_{\mathrm{disc},n}$ is the discrete selection
Rayleigh quotient (finite-stage replication-rate differential)
and $\mathcal{R}_{\mathrm{cts},n}$ is the continuous-time surrogate
from the declared population-interface regime.
\end{definition}

\begin{theorem}[\status{C}]\label{thm:bio-EVO-bridge-template}
\textup{(Evolutionary Bridge Template.)}
\textup{[Conditional on: Thm.~\ref{thm:bio-EVO-discrete-dynamics}
(discrete selection dynamics certified); and the evolutionary
domain loss vanishing: $\varepsilon_n^{\mathrm{evo}} \to 0$
along the hereditary ledger exhaustion.]}
If $\varepsilon_n^{\mathrm{evo}} \to 0$, then the discrete hereditary
transition family $(P_n)_{n\ge0}$ is asymptotically coherent in the
declared population-interface regime, and continuous-time population
dynamics are licensed as effective shorthand by Book~VI.
In particular, the certified finite-stage selection advantage
$\Delta r_n([X],[Y])$ (Def.~\ref{def:bio-differential-replication-advantage})
survives the passage to continuous time without selection-rate loss.
\end{theorem}

\begin{proof}
By Thm.~\ref{thm:bio-EVO-discrete-dynamics}, the discrete selection
dynamics are certified and quotient-visible. The evolutionary domain
loss $\varepsilon_n^{\mathrm{evo}}$ measures how much of the
finite-stage replication-rate differential is lost in the
continuous-time limit. The proof follows the B2-DomainNoLoss
pattern (\texttt{thm:B2-DomainNoLoss}, YM Branch §B2): by the
Rayleigh-characterization of selection dynamics,
$\mathcal{R}_{\mathrm{cts},n}(\psi) \ge
\mathcal{R}_{\mathrm{disc},n}(\psi) - \varepsilon_n^{\mathrm{evo}}$.
Taking infima and the exhaustion limit with
$\varepsilon_n^{\mathrm{evo}} \to 0$ gives selection-rate preservation
in the continuous-time limit. Book~VI licenses the continuous-time
surrogate as effective shorthand for the certified discrete dynamics.
\qed
\end{proof}

\begin{openobligation}[\status{O}]\label{ob:bio-EVO-bridge-cond}
\textup{(O-BIO.EVO-Bridge-$\varepsilon$: Convergence Condition.)}
The remaining Book~VI bridge obligation reduces precisely to:
prove that the evolutionary domain loss satisfies
$\varepsilon_n^{\mathrm{evo}} \to 0$ along the hereditary ledger
exhaustion in the declared population-interface regime.
This is the single named convergence condition blocking the
continuous-time evolutionary dynamics bridge.
\end{openobligation}



%% ---------------------------------------------------------------
\section{Second-Order Multicellular Organization}\label{sec:bio-MULTI-discharge}
%% ---------------------------------------------------------------

\begin{definition}[\status{D}]\label{def:bio-unit}
\textup{(Certified BIO Unit.)}
A \emph{certified BIO unit} is a nondegenerate BIO endpoint object
\[
 U_i \;=\; (X_i,\, \Phi_i,\, \mathcal{M}_i,\, F_i,\, \mathcal{B}_i,\, \mathcal{H}_i)
\]
where $X_i \in \mathcal{C}_{\mathrm{Bio}}^{\mathrm{nd}}$, with
certified replication kernel $\Phi_i$, metabolic ledger $\mathcal{M}_i$,
fidelity margin $F_i > 0$, boundary object
$\mathcal{B}_i = \partial I_{X_i}$, and hereditary descent object
$\mathcal{H}_i$. The term \emph{unit} is not imported biology;
it is the lawful nondegenerate BIO endpoint object.
\end{definition}

\begin{definition}[\status{D}]\label{def:bio-inter-metabolic-ledger}
\textup{(Inter-Unit Metabolic Exchange Ledger.)}
For a finite family $\mathcal{U} = \{U_1,\ldots,U_k\}$ of certified
BIO units, the \emph{inter-unit metabolic exchange ledger}
$\mathcal{M}_{\mathcal{U}}^{\mathrm{inter}}$ consists of certified
resource-transformation entries $m_{ij} : U_i \rightsquigarrow U_j$
that are admissibly testable, boundary-visible, record resource
contributions not internal to either unit alone, are Book~III
transport-accounted, and introduce no-hidden resource channel.
\end{definition}

\begin{definition}[\status{D}]\label{def:bio-inter-hereditary-ledger}
\textup{(Inter-Unit Hereditary Coordination Ledger.)}
The \emph{inter-unit hereditary coordination ledger}
$\mathcal{H}_{\mathcal{U}}^{\mathrm{inter}}$ consists of certified
relations $h_{ij} : \mathcal{H}_i \rightsquigarrow \mathcal{H}_j$
where the hereditary persistence or replication of $U_j$ is affected
by a declared quotient-visible contribution from $U_i$, with each
$h_{ij}$ transport-accounted and no-hidden-channel compliant.
\end{definition}

\begin{definition}[\status{D}]\label{def:bio-coordination-graph}
\textup{(BIO Coordination Graph.)}
The \emph{BIO coordination graph} of $\mathcal{U}$ is
$G_{\mathcal{U}} = (V_{\mathcal{U}},\, E^M_{\mathcal{U}},\, E^H_{\mathcal{U}})$
where $V_{\mathcal{U}} = \{U_1,\ldots,U_k\}$ and the edges record
declared inter-unit metabolic and hereditary dependences.
This is not a primitive topology; it is the finite quotient-visible
record of certified inter-unit exchange.
\end{definition}

\begin{theorem}[\status{C}]\label{thm:bio-inter-metabolic-legitimacy}
\textup{(Inter-Unit Metabolic Legitimacy.)}
\textup{[Conditional on O-BIO.MET discharged and each $U_i$ a
certified BIO unit.]}
$\mathcal{M}_{\mathcal{U}}^{\mathrm{inter}}$ is a lawful second-order
metabolic ledger: each entry is admissible, quotient-visible, and
Book~III-accounted.
\end{theorem}

\begin{proof}
Each $U_i$ is a certified nondegenerate BIO endpoint object; its
boundary $\mathcal{B}_i$ and metabolic ledger $\mathcal{M}_i$ are
source-descended. An inter-unit entry $m_{ij}$ is the same kind of
object as a first-order metabolic entry, with endpoints that are
certified BIO units. Since both endpoints are quotient-visible and
boundary-visible, the entry is observable. Since transport and defect
costs are ledger-recorded, no-hidden resource flow is supplemented.
\qed
\end{proof}

\begin{theorem}[\status{C}]\label{thm:bio-inter-hereditary-legitimacy}
\textup{(Inter-Unit Hereditary Legitimacy.)}
\textup{[Conditional on O-BIO.HER discharged and each $U_i$ a
certified BIO unit.]}
$\mathcal{H}_{\mathcal{U}}^{\mathrm{inter}}$ is a lawful second-order
hereditary coordination ledger.
\end{theorem}

\begin{proof}
At first order, $\mathcal{H}_i$ is lawful by O-BIO.HER: it is obtained
by restricting Book~III transport to the replication-visible sector
with no hidden hereditary channel. An inter-unit hereditary entry
$h_{ij}$ records a dependence between two already-certified hereditary
objects. If the dependence is declared, quotient-visible, and
transport-accounted, it is downstream of the first-order hereditary
ledgers. Any hidden inter-unit hereditary channel would violate the
exhaustive source-image condition (Book~IV) applied to $\mathbf{U}$.
\qed
\end{proof}

\begin{theorem}[\status{C}]\label{thm:bio-multi-boundary}
\textup{(Second-Order Boundary Criterion.)}
\textup{[Conditional on Thm.~\ref{thm:bio-BND-nd} and
$\mathcal{U}$ satisfying the nondegenerate boundary-forcing conditions
at the assembly level.]}
Define the outer support $I_{\mathcal{U}}$ as the minimal support
preserving the joint replication-heredity signature of $\mathcal{U}$.
If the assembly-level metabolic exchange crosses
$\partial I_{\mathcal{U}}$, then the second-order boundary object
\[
 \mathcal{B}_{\mathrm{Bio}}^{(2)}(\mathcal{U}) \;:=\; \partial I_{\mathcal{U}}
\]
is a quotient-visible certified interface.
\end{theorem}

\begin{proof}
The proof repeats the argument of Thm.~\ref{thm:bio-BND-nd} with $X$
replaced by $\mathcal{U}$: the joint replication-heredity signature
plays the role of the first-order replication signature, and the
inter-unit exchange ledger $\mathcal{M}_{\mathcal{U}}^{\mathrm{inter}}$
plays the role of the metabolic ledger. The same Closed-Core Alternative
argument excludes all degenerate cases, giving $\partial I_{\mathcal{U}}$
as a certified interface with Book~II capacity control.
\qed
\end{proof}

\begin{theorem}[\status{C}]\label{thm:bio-multi-endpoint}
\textup{(Conditional Multicellular Endpoint.)}
\textup{[Conditional on Thms.~\ref{thm:bio-inter-metabolic-legitimacy},
\ref{thm:bio-inter-hereditary-legitimacy}, \ref{thm:bio-multi-boundary};
nondegenerate coordination of $\mathcal{U}$; and no-hidden inter-unit
channel.]}
A finite family $\mathcal{U}$ of certified BIO units with lawful
inter-unit metabolic and hereditary ledgers and a second-order boundary
object admits a lawful multicellular BIO endpoint
\[
 E_{\mathrm{Bio}}^{\mathrm{multi}} \;=\;
 \bigl(\mathcal{U},\;
 \mathcal{M}_{\mathcal{U}}^{\mathrm{inter}},\;
 \mathcal{H}_{\mathcal{U}}^{\mathrm{inter}},\;
 \mathcal{B}_{\mathrm{Bio}}^{(2)},\;
 G_{\mathcal{U}}\bigr).
\]
\end{theorem}

\begin{proof}
Certified BIO units supply first-order replication-heredity endpoint.
Inter-unit metabolic and hereditary ledgers supply second-order
coordination. The second-order boundary object is certified by
Thm.~\ref{thm:bio-multi-boundary}. The coordination graph is
finite and quotient-visible. No organismal essence, developmental
program, tissue type, or biological hierarchy is imported.
Book~V branch legitimacy is satisfied: declared endpoint, observable
family visibility, source descent from first-order BIO units,
no-hidden supplement, and no path multiplication.
\qed
\end{proof}

\begin{openobligation}[\status{O}]\label{ob:bio-MULTI-coordination}
\textup{(O-BIO.MULTI-Coord: Inter-Unit Coordination Construction --- Remains Open.)}
The multicellular endpoint theorem requires nontrivial inter-unit
metabolic and hereditary coordination ledgers. Constructing those
ledgers for a specific declared biological target regime remains an
open obligation. The structural architecture is now established
(second-order boundary, inter-unit ledger definitions, coordination
graph), but the coordination must be demonstrated for a concrete
declared assembly.
\end{openobligation}

%% ---------------------------------------------------------------
\section{Updated Closure Ledger}\label{sec:bio-closure-updated}
%% ---------------------------------------------------------------

\begin{remark}[\status{R}]\label{rem:bio-closure-updated}
\textbf{Updated discharge status after synthesis pass.}
After the promotions in this section, the BIO obligation ledger
is updated as follows.

\medskip
\noindent
\begin{tabular}{|l|l|p{7.5cm}|}
\hline
\textbf{Label} & \textbf{Status} & \textbf{Description} \\
\hline
O-BIO.REP & [C] & Replication Certification \\
O-BIO.MET & [C] & Metabolic Bookkeeping \\
O-BIO.FID & [C] & Replication Fidelity Margin \\
O-BIO.NC & [C] & no-hidden Hereditary Channel \\
O-BIO.HER & [C] & Hereditary Descent \\
O-BIO.BND$|_{\mathcal{C}^{\mathrm{nd}}}$ & [C] & Boundary Self-Organization on nondegenerate regime \\
O-BIO.END$^{\mathrm{nd}}$ & [C] & Nondegenerate Biology Endpoint \\
O-BIO.BND-Force & [O] & Boundary forcing for all nontrivial replicators \\
O-BIO.EVO$_{\mathrm{disc}}$ & [C] & Discrete selection dynamics, fitness as rate \\
O-BIO.EVO-Bridge & [O] & Book~VI bridge for continuous-time dynamics \\
O-BIO.MULTI$|_{\mathcal{U}}$ & [C] & Second-order endpoint for declared assemblies \\
O-BIO.MULTI-Coord & [O] & Inter-unit coordination ledger construction \\
\hline
\end{tabular}

\medskip
The full BIO endpoint $E_{\mathrm{Bio}}$ is now certified on the
declared nondegenerate regime $\mathcal{C}_{\mathrm{Bio}}^{\mathrm{nd}}$.
The scope restriction is the honest statement: this is replication-heredity
closure on configurations where replication, metabolism, and external
resource exchange are all genuinely present.
\end{remark}

%% ---------------------------------------------------------------
\section{Developmental Differentiation via SCC Role Carriers}
\label{sec:bio-DEV}
%% ---------------------------------------------------------------

This section establishes the framework for developmental differentiation
as stable role specialization in a coordinated BIO assembly, carried
by minimal SCC-style support rather than imported cell types. No
primitive ``cell type,'' ``developmental program,'' ``organ,''
``tissue,'' or ``fate map'' is assumed. A role is lawful only when
it is visible in the inter-unit ledgers.

\subsection{Developmental Roles}\label{sec:bio-DEV-roles}

\begin{definition}[\status{D}]\label{def:bio-developmental-role}
\textup{(BIO Developmental Role.)}
Let $\mathcal{U} = \{U_1,\ldots,U_k\}$ be a coordinated BIO assembly
with inter-unit ledgers
$\mathcal{M}_{\mathcal{U}}^{\mathrm{inter}}$,
$\mathcal{H}_{\mathcal{U}}^{\mathrm{inter}}$,
coordination graph $G_{\mathcal{U}}$, and second-order boundary
$\mathcal{B}_{\mathrm{Bio}}^{(2)}$.
The \emph{developmental role} of a unit $U_i$ is its
quotient-visible contribution profile:
\[
 \rho(U_i) \;:=\;
 \bigl(
 \mathcal{M}_{\mathcal{U}}^{\mathrm{inter}}|_{U_i},\;
 \mathcal{H}_{\mathcal{U}}^{\mathrm{inter}}|_{U_i},\;
 \mathcal{B}_{\mathrm{Bio}}^{(2)}|_{U_i},\;
 G_{\mathcal{U}}|_{U_i}
 \bigr).
\]
Two units have the same role iff their contribution profiles are
equivalent under declared unit-level probes and transport bookkeeping.
This is intentionally operational: a role is what a unit does in
the certified assembly ledgers, not what a biologist labels it.
\end{definition}

\begin{definition}[\status{D}]\label{def:bio-role-differentiation}
\textup{(Role Differentiation.)}
A coordinated BIO assembly $\mathcal{U}$ exhibits \emph{role
differentiation} when there exist $U_i, U_j \in \mathcal{U}$ with
$\rho(U_i) \not\sim \rho(U_j)$ under declared unit-level probes, and
the distinction persists under declared inter-unit transport and
hereditary descent.
\end{definition}

\begin{definition}[\status{D}]\label{def:bio-scc-role-carrier}
\textup{(SCC Role Carrier.)}
For a role $\rho$, its \emph{SCC role carrier} $K_\rho$ is the
lawful faithful support carrier for the role observable $\tau_{\mathrm{BioRole}}(\rho)$,
whose support consists of:
\begin{enumerate}[label=\textup{(\arabic*)}]
 \item role witness/outcome support;
 \item normalized role-type support;
 \item transported persistence and defect-history support;
 \item ORA/CTM/TIN/ledger compatibility support.
\end{enumerate}
This is the SCC support-sufficiency template specialized to BIO roles.
No primitive cell-type label may appear as load-bearing support.
\end{definition}

\begin{definition}[\status{D}]\label{def:bio-dev-trajectory}
\textup{(Developmental Differentiation Trajectory.)}
A \emph{developmental differentiation trajectory} is a finite sequence
$\mathcal{U}_0 \to \mathcal{U}_1 \to \cdots \to \mathcal{U}_n$ of
coordinated BIO assemblies such that each transition is recorded by
declared inter-unit metabolic and hereditary ledgers, role profiles
$\rho_t(U_i)$ are quotient-visible at each stage, SCC role carriers
$K_{\rho_t}$ persist or lawfully change under transport, and no
primitive fate map or developmental program is imported.
\end{definition}

\subsection{Role Theorems}\label{sec:bio-DEV-theorems}

\begin{theorem}[\status{C}]\label{thm:bio-role-visibility}
\textup{(Role Visibility.)}
\textup{[Conditional on Thms.~\ref{thm:bio-inter-metabolic-legitimacy}
and \ref{thm:bio-inter-hereditary-legitimacy}.]}
If a unit's role is defined by its contribution profile
$\rho(U_i)$, then role identity and role distinction are
quotient-visible.
\end{theorem}

\begin{proof}
Each component of $\rho(U_i)$ is already declared BIO endpoint data
at first or second order: inter-unit metabolism, inter-unit heredity,
second-order boundary contribution, and coordination graph incidence.
The coordination graph is the finite record of certified metabolic
and hereditary edges, not a primitive topology. Therefore role
distinction is witnessed by declared probes over declared ledgers.
If no declared probe separates two contribution profiles, no role
distinction may be drawn.
\qed
\end{proof}

\begin{theorem}[\status{C}]\label{thm:bio-no-celltype-import}
\textup{(No Primitive Cell-Type Import.)}
\textup{[Conditional on Thm.~\ref{thm:bio-role-visibility} and
Sr.~\ref{sr:bio-no-smuggling}.]}
No developmental role in a coordinated BIO assembly may depend on
a primitive cell type, fate label, tissue label, or organismal
program not visible in the contribution profile $\rho(U_i)$.
\end{theorem}

\begin{proof}
Suppose two units agree on all role profile components but are assigned
different developmental roles. The difference is not sourced by
inter-unit metabolism, inter-unit heredity, boundary contribution, or
coordination graph data. It is therefore a hidden target-side
biological label. That violates the BIO no-hidden-channel discipline
(O-BIO.NC, discharged) and the Book~V branch-legitimacy rule excluding
hidden supplementation.
\qed
\end{proof}

\begin{theorem}[\status{C}]\label{thm:bio-role-MCS}
\textup{(SCC Role-Carrier Minimality.)}
\textup{[Conditional on SCC Support Sufficiency and SCC-MCS
available for the BIO role observable
$\tau_{\mathrm{BioRole}}(\rho)$
(\texttt{thm:scc-mcs}, SCC Branch).]}
Every developmental role $\rho$ has a unique minimal lawful SCC
role carrier $M_\rho$ up to role-carrier equivalence.
\end{theorem}

\begin{proof}
Treat the role observable as an SCC-style structural predicate
asking whether a unit's contribution profile realizes a stable
role across assembly transport. The required support is exactly
the four SCC load-bearing components
(Def.~\ref{def:bio-scc-role-carrier}). By SCC Support Sufficiency,
all lawful role support is exhausted by those four classes. By
SCC Support-Rigidity and Shared-Support Faithfulness (SCC Branch),
the common realized-support carrier remains faithful. SCC-MCS then
gives uniqueness of the minimal carrier.
This is essential for anti-plasticity: without minimality, a
``role change'' could be produced by adding hidden support;
with $M_\rho$ in hand, role transitions must be supported by the
unique minimal carrier change.
\qed
\end{proof}

\begin{theorem}[\status{C}]\label{thm:bio-dev-trajectory-legitimacy}
\textup{(Differentiation Trajectory Legitimacy.)}
\textup{[Conditional on Thm.~\ref{thm:bio-role-MCS} and every role
transition $\rho_t(U_i) \to \rho_{t+1}(U_i)$ being recorded by
declared inter-unit ledgers.]}
A developmental differentiation trajectory is lawful if every role
transition is ledger-visible and carried by SCC-minimal role carriers.
\end{theorem}

\begin{proof}
Role transition is visible by Thm.~\ref{thm:bio-role-visibility}.
Hidden cell-type import is blocked by Thm.~\ref{thm:bio-no-celltype-import}.
Minimal role carriers are supplied by Thm.~\ref{thm:bio-role-MCS}.
Since each transition is ledger-visible and transport-accounted, the
trajectory is source-descended and Book~V-compliant.

The SCC contribution is essential: without MCS, a ``role change''
could be produced by adding hidden support outside the four declared
classes. With $M_\rho$ in hand, the role transition must be
supported by the unique minimal carrier change, preventing arbitrary
carrier expansion across stages.
\qed
\end{proof}

\begin{theorem}[\status{C}]\label{thm:bio-dev-endpoint}
\textup{(Conditional Developmental-Differentiation Endpoint.)}
\textup{[Conditional on Thms.~\ref{thm:bio-multi-endpoint},
\ref{thm:bio-role-visibility}, \ref{thm:bio-no-celltype-import},
\ref{thm:bio-role-MCS}, \ref{thm:bio-dev-trajectory-legitimacy};
and the declared role-carrier-stable regime
$\mathcal{C}_{\mathrm{Bio}}^{\mathrm{dev}}$ of coordinated BIO
assemblies with quotient-visible role differentiation and SCC-minimal
role carriers.]}
The developmental endpoint datum
\[
 E_{\mathrm{Bio}}^{\mathrm{dev}} \;=\;
 \bigl(
 E_{\mathrm{Bio}}^{\mathrm{multi}},\;
 \{\rho\},\;
 \{M_\rho\},\;
 \mathcal{T}_{\mathrm{dev}}
 \bigr)
\]
is a lawful developmental-differentiation endpoint: developmental
differentiation is certified as stable role specialization under
certified coordination, not as imported developmental biology.
\end{theorem}

\begin{proof}
The multicellular endpoint $E_{\mathrm{Bio}}^{\mathrm{multi}}$
supplies certified units, inter-unit coordination, and second-order
boundary. Thm.~\ref{thm:bio-role-visibility} supplies role
visibility. Thm.~\ref{thm:bio-no-celltype-import} blocks imported
cell types. Thm.~\ref{thm:bio-role-MCS} supplies SCC-minimal
anti-plasticity for roles. Thm.~\ref{thm:bio-dev-trajectory-legitimacy}
supplies lawful role-transition trajectories. No organism-level
primitive is added; everything is built from declared first-order
units, second-order ledgers, and SCC-minimal role carriers.
Book~V branch legitimacy is satisfied: declared endpoint, observable
family visibility, source descent, no-hidden supplement, and no
path multiplication.
\qed
\end{proof}

\begin{remark}[\status{R}]\label{rem:bio-DEV-SCC-role}
\textbf{SCC architecture and anti-plasticity.}
SCC's utility in the developmental setting is specific and structural:
it provides the anti-plasticity mechanism that prevents role
specialization from becoming arbitrary carrier inflation. Without
SCC-minimal role carriers, one could always inflate a ``role'' by
adding hidden support; the SCC branch's MCS theorem says the minimal
faithful carrier is unique up to equivalence, blocking this inflation.
The synthesis document is explicit on this point: SCC does not act
as a consciousness or endpoint-label mechanism here; its sole
structural contribution is anti-plasticity for role carriers.
\end{remark}

\begin{openobligation}[\status{C}]\label{ob:bio-DEV-scc-dep}
\textup{(O-BIO.DEV-SCC: SCC Role-Carrier Dependency --- Conditionally Resolved.)}
\textup{[Cascade from SCC-MCS discharge (Session~XLII,
\texttt{thm:scc-mcs}).]}
Thm.~\ref{thm:bio-role-MCS} requires SCC-MCS for BIO role
observables. O-SCC.MCS has been conditionally discharged
(\texttt{thm:scc-mcs}, SCC Branch) via the SCC-FL/LB/Zorn
existence argument. The SCC-SV, SCC-TL, SCC-FL, SCC-LB
chain is already certified [C] in the SCC branch. The
branch-local verification required here is that BIO role
observables satisfy SCC support sufficiency (four-component
support: realized spectral/coercive, normalized type, transported
defect-history, ORA/CTM/TIN/ledger compatibility).

By Def.~\ref{def:bio-scc-role-carrier}, the BIO role carrier is
defined with exactly those four support components. Support
sufficiency (Thm.~\ref{thm:bio-role-visibility}) is established:
role identity and distinction are quotient-visible.
Finite localizability follows because each role carrier component
is finite at the certified BIO unit scope. Therefore the
SCC-MCS conditions hold for BIO role observables, and
O-BIO.DEV-SCC is conditionally discharged as a cascade from the
SCC-MCS discharge.
\end{openobligation}

\end{document}
